\section*{\hr{Anhang}}

\section{\hr{Fehlerbehebung}}
\begin{table}[H]
	\begin{tabular}{L{0.29\linewidth} L{0.29\linewidth} L{0.29\linewidth}}
		\toprule
		\textbf{\tr{Problem}} & \textbf{\tr{Mögliche Ursachen}} & \textbf{\tr{Lösung}} \\
		\midrule
		\tr{Das LARUS Vario Display startet, jedoch ist das Satellitenpiktogramm rot und die Vario-Zeiger sind fixiert.}
		& \tr{Anschluss an den LARUS-CAN-Port mit einem gekreuzten Rx/Tx-Patchkabel anstelle eines standardmäßigen 1:1-Patchkabels.}
		& \tr{Bitte ersetzen Sie das Patchkabel und verwenden Sie das im Lieferumfang enthaltene Kabel.} \\
		\tr{Das LARUS Vario Display startet, jedoch ist das Satellitenpiktogramm rot und die Vario-Zeiger sind fixiert.}
		& \tr{Das LARUS Vario Display wurde am falschen Stecker angeschlossen (RS232).}
		& \tr{Bitte verbinden Sie die CAN-Anschlüsse.} \\
		\tr{Das Satellitenpiktogramm ist dauernd oder öfters gelb, die Vario- und/oder Windwerte sind nicht plausibel.}
		& \tr{Schlechter GNSS Empfang}
		& \tr{Sorgen Sie dafür, dass die GNSS Antenne ohne (metallische) Abschirmung nach oben plaziert ist.} \\
		\tr{Vario- und/oder Windwerte sind dauernd oder zeitweise nicht plausibel.}
		& \tr{Die LARUS Sensoreinheit wird durch magnetische Einflüsse gestört.}
		& \tr{Platzieren Sie die LARUS Sensoreinheit nicht in der Nähe von (beweglichen) Eisenteilen oder Magneten.} \\
		\tr{Vario- und/oder Windwerte sind nicht plausibel.}
		& \tr{Die Einbaulage der LARUS Sensoreinheit wurde nicht kalibriert.}
		& \tr{Führen Sie die Kalibrierung durch} (\nameref{sensorunit-calibration}). \\
		\bottomrule
	\end{tabular}
	\caption{\hr{Fehlerbehebung}}
\end{table}
\newpage
